\documentclass[10pt,letterpaper]{article}

\usepackage[margin=0.52in]{geometry}
\usepackage[T1]{fontenc}
\usepackage{lmodern}
\usepackage{enumitem}
\usepackage{titlesec}
\usepackage[hidelinks]{hyperref}
\usepackage{xcolor}

\pagestyle{empty}
\setlength{\parindent}{0pt}
\setlength{\parskip}{0pt}
\setlist[itemize]{leftmargin=1.1em,itemsep=1.2pt,topsep=2pt,parsep=0pt,partopsep=0pt}
\titleformat{\section}{\large\bfseries\uppercase}{}{0pt}{}[\vspace{-3pt}\titlerule]
\titlespacing*{\section}{0pt}{6pt}{3pt}

\newcommand{\role}[4]{%
  \textbf{#1} \hfill \textbf{#2}\\[-1pt]
  \textit{#3} \hfill \textit{#4}\vspace{-2pt}
}

\newcommand{\project}[2]{%
  \textbf{#1} \hfill \textit{#2}\vspace{-2pt}
}

\begin{document}

\begin{center}
  {\LARGE\bfseries RAYYAN SHAN}\\[3pt]
  Edmonton, AB $\mid$ Open to Relocation $\mid$
  \href{tel:+17805668534}{(780) 566-8534} $\mid$
  \href{mailto:rayyanshan027@gmail.com}{rayyanshan027@gmail.com}\\[2pt]
  \href{https://linkedin.com/in/rayyan-shan-65a637239}{linkedin.com/in/rayyan-shan-65a637239} $\mid$
  \href{https://github.com/rayyanshan027}{github.com/rayyanshan027} $\mid$
  \href{https://rayyanshan027.github.io/Portfolio}{rayyanshan027.github.io/Portfolio}
\end{center}
\vspace{-7pt}

\section{Education}
\textbf{University of Alberta} \hfill \textbf{Graduated April 2026}\\[-1pt]
Bachelor of Science in Computing Science, Minor in Mathematics \hfill Edmonton, AB

\section{Technical Skills}
\textbf{Languages:} Python, Java, JavaScript, TypeScript, SQL, C\\
\textbf{Frameworks:} PyTorch, FastAPI, Django REST Framework, React, Vite\\
\textbf{ML/Data:} Computer Vision, Semantic Segmentation, U-Net++, OpenCV, NumPy, Pandas, scikit-learn, XGBoost\\
\textbf{Tools:} PostgreSQL, Docker, Linux, Nginx, Git, GitHub, Hetzner Cloud

\section{Experience}
\role{AI Research Software Developer}{January 2026--Present}{Quorum AI, University of Alberta}{Edmonton, AB}
\begin{itemize}
  \item Built and deployed an AI-powered biomedical image-segmentation platform using PyTorch, U-Net++, FastAPI, React, TypeScript, and Docker.
  \item Developed the complete ML pipeline for 200 expert-annotated microscopy images, including preprocessing, augmentation, validation, inference, confidence estimation, and automated cell-metric generation.
  \item Created REST APIs and researcher workflows for image upload, segmentation visualization, TIFF export, and batch ZIP downloads; deployed to a Linux cloud server with Docker Compose and Nginx.
  \item Supported 2--10 oncology researchers processing 600+ images with substantially less manual segmentation; contributing as a co-author to the resulting manuscript.
\end{itemize}

\role{Software Developer}{June 2026--Present}{Omintel}{Remote}
\begin{itemize}
  \item Extracted and standardized 300+ enterprise AI use cases from technology-partner websites into a structured knowledge base.
  \item Designed a normalized taxonomy and search schema covering industries, business problems, AI capabilities, implementation approaches, and measurable outcomes.
  \item Evaluated LLM-assisted extraction workflows and collaborated with technical and business stakeholders to improve partner onboarding and AI-solution discovery.
\end{itemize}

\section{Projects}
\project{Cross-Session Neural Prediction Pipeline}{Python, scikit-learn, XGBoost}
\begin{itemize}
  \item Built a leave-one-session-out machine-learning pipeline using CRCNS HC-6 neural-recording data and compared Ridge Regression, XGBoost, and nonlinear prediction models.
  \item Developed a residual-adaptation method that reduced error by 25--73\% on difficult sessions and outperformed XGBoost on 8 of 9 evaluated subjects.
\end{itemize}

\project{Cyan Federated Social Platform}{Django REST Framework, React, PostgreSQL}
\begin{itemize}
  \item Developed a federated social platform with authentication, profiles, posts, comments, likes, follows, and server-to-server communication.
  \item Built REST APIs, automated tests, and backend services supporting distributed ActivityPub-style functionality.
\end{itemize}

\section{Publications}
\textbf{AI-based Biomedical Image Segmentation for Automated Analysis of Fluorescence Microscopy Images}\\[-1pt]
Co-author \hfill \textit{Manuscript in preparation}

\end{document}
